\section{Convergent validity and specification choice}
\label{sec:validity}

\S\ref{sec:original} established that \citet{lagasio2024} states neither how per-term
TF-IDF values are aggregated into a document-level substance score nor whether the term
vector is normalised, and \S\ref{sec:measurement} that the second of those silences is
first-order: it determines what the resulting quantity measures. Three substance
specifications therefore remain live and no reading of the original's text selects among
them. This section selects on evidence, using a criterion external to the specifications
compared --- external, but not independent of them, and the two respects in which it is
not are stated below with their magnitudes rather than deferred to
\S\ref{sec:limitations}. One of the two is then removed by recomputation on bases the two
instruments do not share (\S\ref{sec:validity-disjoint}); the other is not removable within
this design and is reported as a bound.

\subsection{\QUANT{} as a yardstick, and the senses in which it is not independent}

\QUANT{} (\S\ref{sec:index}) is a density of regular-expression matches over the
\emph{raw} extracted text: percentage expressions, integers of four or more digits with
calendar years excluded, physical units of mass, energy and volume, and reporting-framework
acronyms, divided by the raw token count $R_d$. Its construction shares almost nothing with
any \SUS{} specification. There is no vectoriser, no lemmatised vocabulary, no
document-frequency weighting and no vector space; the token base is the raw text with
digits intact, whereas every \SUS{} specification is computed on the lemmatised,
stopword-filtered text (\S\ref{sec:lexicon}).

It also targets a property that substantive disclosure should exhibit: a report that
discloses substance quantifies it, in tonnes, percentages and absolute counts, and names
the frameworks under which those figures are computed. Agreement with \QUANT{} is
therefore a convergent-validity test in the sense of \citet{campbell1959} --- two
instruments built by largely unrelated means, aimed at overlapping constructs, should
covary. Agreement does not prove validity, since \QUANT{} is itself a crude instrument and
two measures can covary for irrelevant reasons; but a substance measure that does not
covary with independently counted quantitative content is hard to defend.

``Almost nothing'' is doing work in that claim, and we discharge it before the result rather
than after it. Independence of construction fails in one respect --- fourteen vocabulary
entries are counted by both instruments --- and independence of sampling in another: both
are computed over the same extracted text, with denominators that are two measures of the
same document's length. Both are quantified below, and neither is left to
\S\ref{sec:limitations}.

\begin{table}[htbp]
\centering
\caption{Correlation of each substance specification with \QUANT{} ($n = 343$).}
\label{tab:validity-quant}
\begin{tabular}{@{}lrr@{}}
\toprule
Substance specification & Pearson $r$ & Spearman $\rho$ \\
\midrule
$\sustfidf$ & 0.6817 & 0.7220 \\
$\susdens$ & 0.6663 & 0.7076 \\
$\suslag$ & 0.3217 & 0.4299 \\
Sublinear TF, length-normalised & $+0.0067$ & --- \\
Sublinear TF, mean over vocabulary & $+0.5262$ & --- \\
\bottomrule
\end{tabular}
\end{table}

The two length-normalised specifications track independently measured quantitative
content roughly twice as well as the L2-normalised specification does. The ordering is
unchanged under rank correlation, so it is not produced by a handful of extreme documents,
though the gap narrows there (0.7076 against 0.4299). This is the predicted consequence of
the invariance identity in \S\ref{sec:measurement}: $\suslag$ discards the magnitude of a
document's ESG content and retains only its direction across the vocabulary, and quantified
content is a magnitude.

\subsection{The yardstick is partly contaminated}

The specification choice in this paper rests on agreement with \QUANT{}, and \QUANT{} is not
clean: it counts some of the same strings that the \SUS{} vocabulary counts. The instrument
doing the validating and the instrument being validated share terms. We give the size of the
overlap, the reason the obvious defence of it fails, and then, in
\S\ref{sec:validity-disjoint}, the computation that settles what the defence cannot.

The \texttt{frameworks} group of \QUANT{} matches 11 entries of the 154-term ESG vocabulary
(\emph{csrd}, \emph{gri}, \emph{issb}, \emph{sasb}, \emph{sdg}, \emph{sfdr}, \emph{tcfd},
and four members of the \emph{taxonomy} family); the \texttt{units} group matches three more
(\emph{co2}, \emph{co2e}, \emph{ghg}). Those 14 shared entries account for 67{,}194 of
895{,}996 vocabulary occurrences over the deduplicated corpus, i.e.\ \textbf{7.50\,\%} of
total \SUS{} occurrence mass.
The exposure on the \QUANT{} side we can only bound: the \texttt{frameworks} group, which
holds 11 of the 14 shared entries but also patterns with no vocabulary counterpart
(\emph{un global compact}, \emph{paris agreement}), supplies 17.22\,\% of all \QUANT{} hits
pooled over the corpus, and the \texttt{units} group adds the remaining three. We do not
report the shared entries' exact share of \QUANT{} hits, and it should not be read off the
group totals. What the group totals do show is that the overlap is not stationary: the
\texttt{frameworks} share rises from 10.15\,\% of \QUANT{} hits in 2018 to 21.87\,\% in
2024, so the shared component grows across the study window and must be carried into
\S\ref{sec:temporal}.

There is a tempting defence, and it does not hold. It runs: the 14 entries belong to the
vocabulary, every specification is computed over that same vocabulary, so the overlap
inflates every row of Table~\ref{tab:validity-quant} alike and cannot generate the ordering
between rows. Two things are wrong with it. First, 7.50\,\% is an \emph{occurrence-mass}
share, which is the weight those entries carry under uniform weighting --- that is, under
$\susdens$ and under no other specification in the table. IDF weighting penalises terms
present in nearly every document, and framework acronyms, \emph{co2} and \emph{ghg} are
exactly such terms, so the shared entries carry less than their mass share under
$\sustfidf$ and $\suslag$; and under L2 scaling a term present in almost every document
contributes little cross-sectional variance to a score that has been stripped of magnitude.
Second, and more basic: what displaces a correlation is the covariance the shared terms
carry \emph{within each specification}, and a share of occurrences bounds no covariance. No
arithmetic converts 7.50\,\% of mentions into a maximum displacement of $r$, so the claim
that a shared component of that size could not manufacture the distance between 0.6663 and
0.3217 would be an assertion and not a calculation.

The plausible direction is against us. Both mechanisms above reduce the shared entries'
leverage under $\suslag$ relative to $\susdens$, so if the overlap displaces the table at
all it plausibly lifts the top rows further than the bottom row --- inflating the very gap
this section reports. Every correlation in Table~\ref{tab:validity-quant} is therefore an
upper bound on the level of convergent validity, and we do not defend the ranking by
argument.

\subsection{The table recomputed on disjoint bases}
\label{sec:validity-disjoint}

One computation settles what argument cannot, and we ran it. We deleted the 14 shared
entries from the ESG vocabulary, leaving 140 terms, and restricted \QUANT{} to its two
purely numeric families, \texttt{percentages} and \texttt{large\_numbers}, dropping
\texttt{frameworks} and \texttt{units} in their entirety. The two instruments then share no
string. The restricted \QUANT{} correlates 0.8738 with the full one, so it is a
substantially different criterion and not a cosmetic variant of the same one.
Table~\ref{tab:validity-disjoint} recomputes the comparison on those disjoint bases.

\begin{table}[htbp]
\centering
\caption{Correlation with \QUANT{} on the contaminated and on the fully disjoint bases
($n = 343$). Disjoint bases: ESG vocabulary reduced from 154 to 140 entries, \QUANT{}
restricted to \texttt{percentages} and \texttt{large\_numbers}. Pearson $r$ throughout.}
\label{tab:validity-disjoint}
\begin{tabular}{@{}lrr@{}}
\toprule
Substance specification & Full bases & Disjoint bases \\
\midrule
$\sustfidf$ & 0.6817 & 0.4125 \\
$\susdens$ & 0.6663 & 0.4112 \\
$\suslag$ & 0.3217 & 0.1799 \\
\addlinespace
Ratio, $\susdens$ to $\suslag$ & 2.07$\times$ & 2.29$\times$ \\
\bottomrule
\end{tabular}
\end{table}

Read it precisely, because it concedes one thing and establishes another. Every absolute
correlation falls, by roughly two fifths --- so the overlap does inflate the level of
agreement, the figures in Table~\ref{tab:validity-quant} are upper bounds exactly as the
previous subsection said, and the level of convergent validity on clean bases is moderate
rather than strong. What the overlap does \emph{not} do is manufacture the ordering. The
two length-normalised specifications remain roughly twice as close to independently counted
quantitative content as the L2-normalised one, and the relative distance is slightly
\emph{larger} once the shared strings are gone, 2.29$\times$ against 2.07$\times$. That is
the opposite of what the previous subsection expected: on the mechanisms set out there the
overlap should have widened the gap, and in fact it was compressing it slightly. We record
that we had the sign wrong, which is the reason for running the computation rather than
arguing about it. The ranking on which this section's selection rests survives the removal
of everything the two instruments had in common; its level does not, and we do not claim it.

The narrow gap between $\sustfidf$ and $\susdens$ narrows further, to 0.0013, which
reinforces rather than disturbs the reasoning of \S\ref{sec:validity-adopted}: the criterion
separates the length-normalised pair from $\suslag$ and does not separate the pair.

Neither table, in any case, reaches the part of the diagnosis that uses no yardstick at all.
The invariance identity of \S\ref{sec:measurement} is algebraic rather than correlational:
$\suslag$ returns the same value for any two documents whose ESG counts are proportional ---
demonstrated end to end at a deviation of $6.9 \times 10^{-18}$ under twentyfold
amplification, and unchanged to ten decimal places under a 3.48-fold dilution that drops
measured density from 3.0785 to 0.8845 per 100 words. And $\suslag$ correlates 0.9733 with
$\BREADTH$, a diversity statistic computed with no regex, no raw text and no \QUANT{}.
Convergent validity here corroborates a diagnosis that was reached without it. The levels in
Table~\ref{tab:validity-quant} are inflated; the ordering is not, and the diagnosis does not
stand on either table alone.

\subsection{Independence of construction is not independence of sampling}

\QUANT{} and \SUS{} are computed over the same extracted zones (\S\ref{sec:extraction}).
Neither instrument sees text the other does not. A report from which extraction recovered
little usable text scores low on both for reasons that have nothing to do with disclosure
quality, and any bias in the extraction stage --- its density threshold, its context window
--- propagates identically to both sides of the comparison. The test establishes convergent
validity conditional on the extraction stage, not agreement with anything outside our
pipeline. We have no external ESG-washing benchmark; neither does the original
(\S\ref{sec:limitations}).

The denominators sharpen the same point. \QUANT{} is a count over raw tokens $R_d$;
$\susdens$ and $\sustfidf$ are counts over lemmatised tokens $L_d$. These are not the same
variable, but they are two measures of the length of one document and they vary together,
and two ratios with collinear denominators correlate for that reason alone. The
specifications that win the test are the ones built like the yardstick, and $\suslag$, being
scale-free, cannot benefit from that mechanism even in principle.

The corpus bounds how much of the result this can explain. The fourth row of
Table~\ref{tab:validity-quant} is a specification whose cross-sectional variation is
dominated by its length denominator --- a numerator confined to a narrow range over a length
that is not, as the next subsection establishes --- and it correlates $+0.0067$ with
\QUANT{}. Were \QUANT{}'s variation across
documents driven by its own length denominator, that specification would correlate strongly
with it. It does not, which puts the common-denominator artefact near zero on this corpus.
That is an argument from the data, not a formal partialling-out of length, which we do not
report.

\subsection{Why sublinear TF fails, and where}

Raw density invites an obvious objection: a report that repeats one commitment ten times
is scored as though it disclosed ten times as much. The standard remedy is sublinear
term-frequency scaling \citep{manning2008}, which replaces the count $c_{dj}$ by
$1 + \log c_{dj}$ so that the eleventh mention contributes far less than the second.
Applied with the same length denominator,
\begin{equation}
\label{eq:sublinear}
\SUS^{\mathrm{sub}}_d \;=\; \frac{100}{L_d}\sum_{j:\,c_{dj}>0}\bigl(1+\log c_{dj}\bigr),
\end{equation}
the result is not merely weaker but null: $r = +0.0067$ with \QUANT{}, indistinguishable
from zero and two orders of magnitude below $\susdens$.

The failure has to be localised before anything is concluded from it, because two distinct
quantities are involved and they behave in opposite directions. The vocabulary is curated at
154 terms and a document contains on average 87.75 of them, so the sum in the numerator of
\eqref{eq:sublinear} ranges over a bounded and comparatively narrow interval: log damping
compresses each term's contribution towards its mere presence, and the sum accordingly
tracks the number of \emph{distinct} vocabulary terms in the document rather than the volume
of mentions. Document length is under no such bound. The sign change is the diagnostic.
Taken before the division, the numerator of \eqref{eq:sublinear} moves with the
distinct-term count; taken after it, the finished score correlates $-0.3214$ with that same
count, because variation in $L_d$ now dominates the ratio. The two figures describe
different quantities and neither is interpretable without saying which. What a user of the
specification would compute is the second: a measure that falls with document length
irrespective of content.

The defect is therefore the interaction and not the damping. Aggregated as a mean over the
vocabulary instead of divided by document length, sublinear TF correlates $+0.5262$ with
\QUANT{} --- the same damping, a different denominator, most of the agreement restored.
Sublinear weighting is unusable in this configuration; it is not unusable. We report the
distinction because the null result, stated at the level of a method family, would be false.

One kind of evidence is excluded here by design. The specifications also differ in the shape
of the temporal series they produce, and that behaviour is used neither for nor against any
of them in this section. A specification selected because it delivered the trend the paper
later reports would make that trend an artefact of the selection.
\S\ref{sec:temporal} reports the series and \S\ref{sec:robustness} the sensitivity of the
index to the specification; the selection uses Table~\ref{tab:validity-quant} and
\S\ref{sec:measurement} only.

\subsection{The repetition objection, and what would answer it}

The objection that motivated sublinear TF is not thereby answered --- it is misdirected.
Damping treats repetition as measurement noise to be suppressed inside the substance
score. But a report that restates the same commitment across five sections is exhibiting
a behaviour of interest, plausibly correlated with the construct the index is meant to
detect, and folding it into the substance measure discards that information. The
appropriate response is to measure redundancy as a quantity in its own right --- through
the concentration of the term distribution, or near-duplicate detection across sections
--- and to enter it in the index alongside substance. $\BREADTH$ (\S\ref{sec:index}) is
a first step in that direction and is the reason we report it, but it is reported rather
than incorporated. We implement no redundancy component and claim no result about one;
this is recorded as the direction the objection actually points.

\subsection{Specification adopted}
\label{sec:validity-adopted}

Table~\ref{tab:validity-quant} ranks $\sustfidf$ first, at 0.6817 against 0.6663. We adopt
$\susdens$ nonetheless, and the reason is not that the difference is negligible. The gap is
0.0154 between two measures that correlate 0.9926 with each other; whether it is
statistically distinguishable is a question about dependent correlations that we do not
test, so we neither claim the two specifications are equivalent on the criterion nor treat
the ordering between them as informative. The criterion's discriminating power in this
section lies between the length-normalised pair and $\suslag$ --- a factor of two --- and
not within the pair.

The tie is therefore broken on a structural ground rather than an evidential one. $\susdens$
estimates nothing on the sample: it requires no document frequencies, so a report's score
does not change when the corpus grows and is comparable across corpora, whereas both
IDF-weighted specifications are defined relative to these 343 reports
(\S\ref{sec:limitations}). Between two specifications the convergent test does not separate,
we take the one that does not carry the sample around with it.

Little rests on that choice and a great deal rests on the other one. The two
length-normalised substance scores correlate 0.9926 and the labels they produce differ on 6
of 343 documents, figures \S\ref{sec:robustness} establishes; $\suslag$ differs from the
adopted specification on 96 of 343. \S\ref{sec:robustness} quantifies that consequence.
